\documentclass[10pt]{article}
\usepackage{resume_shared}

\begin{document}
\fontsize{9.8pt}{11.7pt}\selectfont

\ResumeName{Sri Ravi Kumar Birada}
\ResumeContact{Hyderabad, India}{(+91) 7200999942}{ravikumar31birada@gmail.com}{www.linkedin.com/in/ravi-20043}

\section{Summary}
Enterprise AI solution architect with 7+ years translating product, customer, and platform requirements into scalable AI solution patterns. Strong at shaping architecture narratives across ideation, MVP, evaluation, and scale; aligning AI workflows with integration, governance, and security needs; and converting ambiguous client problems into implementation-ready blueprints. Experience includes agentic AI, RAG, partner readiness, AI roadmap definition, and secure enterprise deployment positioning.

\section{Professional Experience}
\RoleEntry{AI Product Manager - Enterprise AI Products}{HCLTech}{2023 -- Present}
\begin{itemize}
  \item Translated 20+ mapped workflow activities into an AI architecture roadmap across 3 enterprise SaaS product modules, covering 5 use cases over a 12-month planning horizon.
  \item Engineered an agentic AI RFP response workflow that ingests customer RFPs and enterprise capability / technology-stack data, orchestrates retrieval, and applies validation plus sanity checks before final document submission.
  \item Modeled a 0 to 1 enterprise SaaS capability by consolidating inputs from 6+ customer accounts and defining 8 MVP feature areas across platform-extension and standalone solution patterns.
  \item Assessed 25+ GenAI partners and onboarded 3 through technical-fit reviews, 3 pilot-account selections, integration readiness checks, and 5 pilot success criteria.
  \item Designed AI output-control patterns across retrieval relevance, document completeness, hallucination risk, human review, and final artifact readiness for enterprise workflows.
  \item Embedded governance-ready design considerations such as fallback logic, access controls, validation checkpoints, human oversight, and deployment constraints into AI-enabled solution narratives.
\end{itemize}

\RoleEntry{Product Manager - Enterprise SaaS Products}{HCLTech}{May 2022 -- 2023}
\begin{itemize}
  \item Consolidated customer-account requirements into modular enterprise SaaS patterns across customer operations, workflow orchestration, reporting, and commercial governance.
  \item Positioned multi-tenant SaaS capabilities across cloud, API-led integration, reporting, and modular packaging, helping enterprise stakeholders understand implementation paths and adoption trade-offs.
  \item Supported rapid rollout for an enterprise product deployment supporting 300K+ SKUs and 52+ distributed business entities by aligning execution dependencies and reporting enablement, resulting in launch within 21 days.
\end{itemize}

\MiniHeading{Selected Architecture Outcomes}
\begin{itemize}
  \item AI architecture roadmap: aligned 5 AI use cases across 3 enterprise SaaS modules with data, integration, security, and operating-model dependencies.
  \item Agentic RFP workflow: combined document ingestion, enterprise knowledge retrieval, validation gates, and final-output sanity review into a practical proposal automation pattern.
  \item 0 to 1 modular architecture: shaped 8 MVP feature areas from 6+ customer-account inputs into a reusable SaaS capability pattern.
  \item Partner readiness: converted 25+ partner evaluations into 3 onboarded partners, pilot-account selection, integration-fit checks, and 5 success criteria.
\end{itemize}

\RoleEntry{Network Engineer}{Ericsson Global India Limited}{Oct 2016 -- Jan 2019}
\begin{itemize}
  \item Reduced deliverable turnaround time by 15\% by automating report validation workflows with Python.
  \item Cut server infrastructure requirements by 40\% through performance improvements in analysis tooling and workflow design.
  \item Received the Ericsson Power Award in Q4 2017 for automation impact in Single Site Verification reporting.
\end{itemize}

\section{Education}
\InlineSectionText{MBA, MDI Gurgaon (2022)}
\InlineSectionText{B.Tech, Electronics and Communication Engineering, Amrita School of Engineering (2016)}

\section{Certifications}
\InlineSectionText{OpenAI AI Technical Practitioner (2026) \textbar\ Google Certified GenAI Leader (2025) \textbar\ GenAI Prompt Engineer L2 (2025) \textbar\ Certified Scrum Product Owner (CSPO) (2024)}

\section{Skills}
\SkillGroup{Architecture and Design}{AI Solution Architecture, Enterprise Architecture, Systems Integration, Cloud Architecture, Workflow Orchestration, Technical Discovery, Implementation Readiness, Modular Architecture}
\SkillGroup{Enterprise AI}{RAG, Agentic Workflow Design, Knowledge Search, Governance and Controls, Responsible AI, LangChain, CrewAI}
\SkillGroup{Platforms and Tools}{Azure, AWS Positioning, API-led Integration Framing, Partner Onboarding, Pilot Success Definition, Value Stream Mapping, Power BI, Python}

\end{document}
